% ============================================================================================
\section{S2: Contour Mapping (Owner: Shrey Modi, BT2024125)}
\label{sec:s2}
\subsection{Choice of the scalar field}
We contour the daily mean \textbf{vapour pressure deficit (VPD)}. VPD is the difference between saturation and actual vapour
pressure, so it measures how strongly the air pulls water out of plants and fuels~\cite{novick2016,williams2019}. We screened
the three candidates from Sec.~\ref{sec:intro} (Fig.~\ref{fig:10}, Table~\ref{tab:s2choice}).
\begin{itemize}
\item \textbf{rmax} saturates: 21.3\,\% of cell-days equal exactly 100\,\%. Its isolines at the upper deciles fragment into
  118,998 segments along the edges of the plateau, which is noise rather than structure.
\item \textbf{srad} is smooth (roughness 0.121), but its pattern is transient cloud cover with a narrow dynamic range
  (p95/p5 = 2.6).
\item \textbf{VPD} has the largest dynamic range (p95/p5 = 11.0), no plateau and a strong, physically meaningful
  west--east gradient. It also has literature thresholds for plant stress, and it ties S1 (temperature) to moisture.
\end{itemize}

\figfull{Fig_10_s2_variable_candidates.png}{Screening of candidate contour variables on D2. Colour = field (1st--99th
percentile stretch); black lines = our Marching Squares at the 10/30/50/70/90th percentile of the day (also marked on the
colour bars).}{fig:10}
\input{generated/tab_s2_choice}

\subsection{Marching Squares implementation}
\fpath{scivis/s2_contours/marching_squares.py} implements the algorithm from scratch in NumPy, the 2-D analogue of marching
cubes~\cite{lorensen1987}. It is fully vectorised over all $584\times1385$ cells per isovalue (Algorithm~\ref{alg:ms}).
\begin{itemize}
\item \textbf{Case index.} Corners are bottom-left $z_{bl}=Z_{i,j}$, bottom-right $z_{br}=Z_{i,j+1}$, top-right
  $z_{tr}=Z_{i+1,j+1}$ and top-left $z_{tl}=Z_{i+1,j}$. A corner is ``inside'' if $z\ge c$, and the cell's case is
  $\kappa=1[z_{bl}{\ge}c]+2[z_{br}{\ge}c]+4[z_{tr}{\ge}c]+8[z_{tl}{\ge}c]\in\{0..15\}$ (Fig.~\ref{fig:11}).
\item \textbf{Lookup table.} Cases 0 and 15 produce nothing. Each remaining unambiguous case pairs two of the edges
  B, R, T, L, and complementary cases share their segment (e.g.\ 1 and 14 both give L--B).
\item \textbf{Interpolation.} On an edge from $(x_a,z_a)$ to $(x_b,z_b)$ the crossing lies at
  \begin{equation}
  x^\ast = x_a + t\,(x_b-x_a),\qquad t=\frac{c-z_a}{z_b-z_a}\in[0,1].
  \end{equation}
\item \textbf{Saddles (cases 5 and 10).} The two inside corners lie on a diagonal, and the four crossings can be paired
  in two topologically different ways. We resolve this with the \emph{asymptotic decider}~\cite{nielson1991}. Inside the
  cell the bilinear interpolant has a saddle value
  \begin{equation}
  s=\frac{z_{bl}z_{tr}-z_{br}z_{tl}}{z_{bl}+z_{tr}-z_{br}-z_{tl}},
  \label{eq:asym}
  \end{equation}
  and if $s\ge c$ the inside corners are \emph{joined} through the cell; otherwise they are separated. For comparison we also
  implemented the cell-centre mean decider, $\bar z\ge c$, and the fixed rules ``always join'' and ``always separate''.
  The fixed rules are what \texttt{skimage} calls \texttt{fully\_connected='high'} and \texttt{'low'}.
\item \textbf{NaN handling.} A cell with any NaN corner is skipped, so isolines stop at the coastline and the mask edge
  and never cross the ocean.
\end{itemize}

\begin{algorithm}[!t]
\caption{Vectorised Marching Squares with asymptotic decider}
\label{alg:ms}
\begin{algorithmic}[1]
\REQUIRE grid $Z\in\mathbb{R}^{n_y\times n_x}$ (NaN = no data), coordinates $x,y$, isovalue $c$
\ENSURE set of segments $\mathcal{S}$
\STATE $Z_{bl},Z_{br},Z_{tr},Z_{tl}\leftarrow$ the four shifted views of $Z$ \COMMENT{shape $(n_y{-}1)\times(n_x{-}1)$}
\STATE $V\leftarrow$ all four corners finite
\STATE $\kappa\leftarrow [Z_{bl}{\ge}c] + 2[Z_{br}{\ge}c] + 4[Z_{tr}{\ge}c] + 8[Z_{tl}{\ge}c]$
\STATE $P_e\leftarrow$ interpolated crossing on each edge $e\in\{B,R,T,L\}$ of every cell
\FORALL{unambiguous case $k$ with segment list $\mathrm{LUT}[k]$}
  \FORALL{$(e_1,e_2)\in\mathrm{LUT}[k]$}
    \STATE $\mathcal{S}\leftarrow\mathcal{S}\cup\{(P_{e_1}[q],P_{e_2}[q]) : \kappa[q]=k \wedge V[q]\}$
  \ENDFOR
\ENDFOR
\FORALL{saddle case $k\in\{5,10\}$}
  \STATE $s\leftarrow$ Eq.~\eqref{eq:asym} on cells with $\kappa=k \wedge V$
  \STATE add $\mathrm{JOIN}[k]$ segments where $s\ge c$ and $\mathrm{SEP}[k]$ segments elsewhere
\ENDFOR
\RETURN $\mathcal{S}$
\end{algorithmic}
\end{algorithm}

\figfullw{Fig_11_marching_squares_cases.png}{The 16-case lookup table. Orange corners are $\ge$ the isovalue. Green
segments join edge crossings (drawn at edge midpoints for clarity; the code interpolates linearly). For the ambiguous
saddles 5 and 10 the dashed purple pair shows the alternative resolution.}{fig:11}{0.62}

\subsection{Validation against reference implementations}
We compared the output with \texttt{skimage.measure.find\_contours}~\cite{vanderwalt2014} (polylines converted to segments)
and with \texttt{contourpy} on the full D2 VPD field, at three isovalues offset by half the 0.01\,kPa storage step so that
no data value equals the isovalue exactly (Fig.~\ref{fig:12}, Table~\ref{tab:msval}).

\figfull{Fig_12_ms_validation.png}{Validation on D2. (a) Our isolines (thick, coloured) under \texttt{skimage}'s (thin
black) at 1.005, 2.005 and 3.005\,kPa; they coincide everywhere. (b, c) A $9\times9$-cell window in Nebraska containing
saddle cells (red dashed squares) at 2.005\,kPa: (b) the fixed ``join'' rule used by \texttt{skimage}; (c) our asymptotic
decider.}{fig:12}

\input{generated/tab_msval}

\textbf{Inferences.}
(1) With the same saddle rule our implementation is \emph{identical} to \texttt{skimage}. All 17,993 / 21,715 / 11,877
segments match exactly at the three levels (100\,\%), with a maximum vertex distance of 0.000 cells, so the lookup table,
the interpolation and the NaN handling are correct.
(2) Saddles are rare but not negligible: 451, 386 and 179 cells, i.e.\ 2.5, 1.8 and 1.5\,\% of the cells crossed by each
isoline. The asymptotic decider joins 195, 177 and 84 of them (43--47\,\%); the centre-mean decider joins 186, 181 and 81. The
two deciders disagree on only a few saddles, because the cell mean is a reasonable estimate of the bilinear saddle value
when the field is smooth at 4\,km.
(3) The saddle choice changes topology, not geometry. Our asymptotic result differs from \texttt{skimage}'s fixed rule in
836--1,024 segments, yet total isoline length differs by only 1.0--1.5\,\%, and contourpy's length lies between the two.
In Fig.~\ref{fig:12}b the fixed rule draws small diamond-shaped closed loops around single cells. The asymptotic decider
(Fig.~\ref{fig:12}c) instead links them into continuous bands that follow the bilinear interpolant. Both are valid
contours of a piecewise-bilinear field, but only the decider is consistent with that interpolant.
(4) Speed: our NumPy version needs 0.08--0.13\,s per level on 810,810 cells, against 0.02--0.05\,s for the Cython
\texttt{skimage}. A 3$\times$ slowdown is acceptable for a pure-Python implementation.

\subsection{Isolines vs.\ filled contours}
The assignment allows Marching Squares and/or contour filling. We use both, and Fig.~\ref{fig:13} shows why. Our isolines
are drawn with Matplotlib \texttt{LineCollection}. Filled bands use \texttt{contourf}, computed on the projected grid
coordinates. Contouring in longitude/latitude and re-projecting made Cartopy invert some polygons with holes, filling the
whole map, so we avoid that path.

\figfull{Fig_13_isolines_vs_filled.png}{VPD on D1 (8~Jul) with ten levels (0.5--5\,kPa). (a) Isolines only (own Marching
Squares, coloured by level); (b) filled bands only (\texttt{contourf}); (c) filled bands plus isolines plus labels at 1--4\,kPa
(own label placement on the longest joined polylines).}{fig:13}

\textbf{Inferences and rationale.}
(1) Isolines encode \emph{gradient}: their spacing is inversely proportional to $|\nabla\mathrm{VPD}|$. In
Fig.~\ref{fig:13}a they pack tightly along the Sierra Nevada and Central Valley margin, the Mogollon Rim and the Wasatch,
where VPD changes by more than 3\,kPa over a few grid cells (topographic control: hot, dry valleys next to cooler
ridges). Isolines also leave the underlying field visible and suit overlays (Fig.~\ref{fig:16}). However, ten coloured
lines are hard to read as values, and in the humid East the 0.5 and 1.0\,kPa lines form a dense tangle of small closed
loops.
(2) Filled bands encode \emph{magnitude} and \emph{area}. Fig.~\ref{fig:13}b shows at once that the Mojave and Sonoran
deserts and the Central Valley exceed 5\,kPa, while almost everything east of about 97\,$^\circ$W is below 1--1.5\,kPa (the Gulf
moisture of D1, the day Beryl came ashore). The shading hides gradient steepness and makes noisy regions look blotchy.
(3) Combining them (Fig.~\ref{fig:13}c) gives magnitude from colour and gradient and exact values from labelled lines, so
this is our default. The rule we follow: use isolines when the task is locating thresholds or gradients, or overlaying a
second field; use filled contours when the task is reading magnitudes or comparing areas.

\subsection{Number and placement of contour levels}
Fig.~\ref{fig:14} and Table~\ref{tab:levels} compare uniform levels (every 0.5\,kPa), quantile levels (deciles of the day)
and domain thresholds (0.5/1/2/3/4\,kPa, which bracket common plant-stress and fire-weather categories).


\figfull{Fig_14_contour_levels.png}{Level placement on D2. (a--c) Filled contours plus our isolines for the three
schemes; (d--f) land-cell histogram with the levels (vertical lines) and the min/max share of land per band.}{fig:14}
\input{generated/tab_s2_levels}

\textbf{Inferences.}
(1) The VPD histogram on D2 is \emph{bimodal}, with peaks at about 0.9\,kPa (humid East) and about 2.2\,kPa (dry West and Plains)
(Fig.~\ref{fig:14}d). Uniform levels therefore put only two levels in the first mode but six in the sparse tail
above 2.5\,kPa. Band shares range from 0.7 to 26.4\,\%.
(2) Quantile levels give every band 9.6--10.2\,\% of land, but five of the nine levels fall between 0.53 and 1.37\,kPa. The
East breaks into many small contours (147,675 segments, the most). Meanwhile everything above 3.16\,kPa (the Great
Basin, the desert Southwest and the Montana heat dome) collapses into one dark class.
(3) Domain thresholds need the fewest segments (69,862) and give named categories. The 2\,kPa line cleanly separates the
semi-arid West and the northern-Plains heat dome from the humid East. The $\ge$4\,kPa band covers 3.5\,\% of land, exactly
the hot deserts. This is the most interpretable choice, and we use 1/2/3/4\,kPa for the time series (Fig.~\ref{fig:17}).

\subsection{Smoothing before contouring}
We apply a \emph{mask-aware} Gaussian filter, $\tilde f=G_\sigma*(f\,m)\,/\,G_\sigma*m$ with $m$ the land mask, so that ocean
NaNs do not bleed into the coast. Then we contour at 2 and 3\,kPa (Fig.~\ref{fig:15}, Table~\ref{tab:smooth}).

\figfull{Fig_15_contour_smoothing.png}{Effect of Gaussian pre-smoothing ($\sigma$ = 0, 1, 3, 6 cells $\approx$ 0, 4, 12,
24\,km) on the 2\,kPa (blue) and 3\,kPa (red) isolines on D2. Top: CONUS (box = zoom); bottom: Utah/Colorado
zoom.}{fig:15}
\input{generated/tab_s2_smooth}

\textbf{Inferences.}
(1) Unsmoothed isolines (33,599 segments) follow every ridge and valley. In the Utah/Colorado zoom they form 311 separate
polylines with about 17,700\,km of total length, most of them small closed loops around individual peaks and basins.
This is real terrain-driven VPD, but it is visual noise at CONUS scale.
(2) $\sigma=1$ cell already cuts the zoom to 58 polylines. $\sigma=3$ (12\,km) keeps the main structure, the dry Colorado
Plateau and Uinta Basin against the moist Colorado Rockies, with 13 lines and 42\,\% of the original length. This is the
best compromise.
(3) $\sigma=6$ (24\,km) leaves 9 lines. Separate dry basins merge, and the 3\,kPa line shifts into the mountains,
misplacing the threshold by tens of kilometres. Smoothing trades positional accuracy for legibility, so the smoothing
scale should be stated whenever smoothed contours are used for decisions.

\subsection{Overlay of contours on colour-mapped fields}
\figfull{Fig_16_contour_overlay.png}{Contours overlaid on colour maps for D1. (a) Labelled VPD isolines (1--4\,kPa) over the
VPD field; (b) the same isolines (coloured by level) over the $\Tmax$ field from S1, giving a bivariate map.}{fig:16}

\textbf{Inferences.}
(1) In Fig.~\ref{fig:16}a the 4\,kPa contour encloses the Mojave and Sonoran deserts and the southern Central Valley. The
1\,kPa contour meanders along about 100\,$^\circ$W and then follows the Appalachians.
(2) The bivariate overlay (Fig.~\ref{fig:16}b) shows that VPD tracks temperature in the West: the $\ge$3--4\,kPa areas sit
exactly on the $\ge$40\,\degC{} areas of California, Nevada and Arizona, as expected because saturation vapour pressure grows
exponentially with temperature. In the East it does not. Texas, Arkansas and the Southeast reach 33--38\,\degC{} on D1, yet
VPD stays below 1--2\,kPa because Beryl's tropical air mass is nearly saturated. ``Hot'' and ``dry'' are therefore
different stresses, and only the overlay shows where they diverge.

\subsection{Contour evolution over the eight dates}
Fig.~\ref{fig:17} uses \emph{fixed} levels (1/2/3/4\,kPa) on every date, which is required for frames to be
comparable.\footnote{Animation: \fpath{images/Fig_17_vpd_contour_evolution.gif} (8 frames, 1 frame/s).}
Table~\ref{tab:evol} gives the area above each level.

\figfull{Fig_17_vpd_contour_evolution.png}{VPD contours on the eight dates with fixed levels (1/2/3/4\,kPa): our isolines
plus filled bands. All eight frames of the GIF \fpath{Fig_17_vpd_contour_evolution.gif}.}{fig:17}
\input{generated/tab_s2_evol}

\textbf{Inferences.}
(1) The share of land with VPD $\ge$2\,kPa peaks on D2 (34.2\,\%), when the northern-Plains heat dome (Montana, Wyoming,
the Dakotas) reaches $\ge$3--4\,kPa, and declines to 16.3\,\% on D7. The extreme $\ge$4\,kPa class shrinks from 4.6\,\% (D1)
to 1.1--2.0\,\% in Aug--Sep.
(2) The $\ge$1\,kPa share barely moves (63.5--72.6\,\%). The arid West is persistently dry, and the season mainly changes
how \emph{intense} the dryness is, not \emph{where} it is.
(3) Tropical systems leave a clear moist footprint. On D3 the Southeast coast falls below 1\,kPa under Debby, and on D7 and D8
the whole Southeast and Appalachians do so under Helene. On D4 and D5 a dry tongue of 2--3\,kPa pushes east from Texas into
Oklahoma and Arkansas, a late-summer dry spell between the storms.
(4) The isoline segment count falls from 57,087 (D2) to 36,559 (D7). The September fields are smoother, with fewer
convective moisture patches.

\subsection{Observations and limitations (S2)}
Our Marching Squares reproduces the reference exactly and adds a principled saddle decider. Filled and line contours
answer different questions, so we combine them. Level placement should follow the data distribution \emph{and} the
domain, and smoothing must be reported because it moves thresholds. Limitations: VPD is a daily mean, so afternoon peaks are
higher than shown. Segments are drawn unjoined in the maps; we join them only for labelling and counting. Contouring on the
geographic grid ignores the slight cell-area change with latitude.
